\chapter{Linear Codes}

Linear codes are one of the most important concepts in classical coding theory. Their efficiency in encoding information and detecting errors makes them computationally highly practical. More importantly, the general concept of linear codes is foundational to quantum error-correcting codes, specifically through their generalization into stabilizer codes. Therefore, it is important to take a brief detour from pure quantum codes to review the definition, characteristics, and advantages of linear codes, establishing the groundwork for how stabilizer codes operate.

\begin{definition}\label{def:linear_code}
  A binary code $\mathcal{C}$ that encodes $k$ bits of information into an $n$-bit space is called a \textit{linear code} if it forms a $k$-dimensional subspace of the vector space $\mathbb{F}^n_2$. Every possible $k$-bit string to be encoded can be represented as a column vector $u \in \mathbb{F}^k_2$, and its corresponding codeword $v \in \mathcal{C}$ is generated by the operation $v = G u$. Here, $G$ is an $n \times k$ matrix with entries in $\mathbb{F}_2$, called the \textit{generator matrix}. The column vectors of $G$ form a basis for the subspace $\mathcal{C}$. Thus, every operation is performed modulo two, and every valid codeword is a linear combination of these basis vectors \citep{Gottesman1998}.
\end{definition}

This definition highlights how linear codes streamline the encoding process by reducing it to matrix multiplication over $\mathbb{F}_2$, an operation that is highly efficient to implement computationally. However, the most profound connection to stabilizer codes lies in the error detection phase, specifically through the use of a parity-check matrix.

\begin{lemma}\label{lemma:complement_of_orthogonal_complement}
Given a subspace $\mathcal{C} \subseteq \mathbb{F}^n_2$ and its orthogonal complement $\mathcal{C}^\perp$, it holds that:
\begin{equation}
(\mathcal{C}^\perp)^\perp = \mathcal{C}
.\end{equation}
\begin{proof}
By the definition of the orthogonal complement, for any $v \in \mathcal{C}$ and $w \in \mathcal{C}^\perp$, their inner product $\langle w, v \rangle = 0$. Therefore, every $v \in \mathcal{C}$ is orthogonal to everything in $\mathcal{C}^\perp$, which implies $\mathcal{C} \subseteq (\mathcal{C}^\perp)^\perp$. 

By calculating dimensions, we observe that:
\begin{equation}
\dim((\mathcal{C}^\perp)^\perp) = \dim(\mathbb{F}^n_2) - \dim(\mathcal{C}^\perp) = n - (n - \dim(\mathcal{C})) = \dim(\mathcal{C})
.\end{equation}
Since $\mathcal{C} \subseteq (\mathcal{C}^\perp)^\perp$ and both subspaces share the same dimension, they must be equal.
\end{proof}
\end{lemma}

\begin{theorem}\label{thm:parity_check_matrix}
Given an $[n, k]$ linear binary code $\mathcal{C}$ with generator matrix $G \in \mathbb{F}_2^{n \times k}$, there exists an $(n - k) \times n$ matrix $H$, called the \textit{parity-check matrix}, such that a vector $v \in \mathbb{F}^n_2$ is a valid codeword ($v \in \mathcal{C}$) if and only if $H v = 0$.
\begin{proof}
Since a linear code $\mathcal{C}$ is a $k$-dimensional subspace of $\mathbb{F}^n_2$, we can consider its orthogonal complement $\mathcal{C}^\perp \subseteq \mathbb{F}^n_2$, defined as:
\begin{equation}
\mathcal{C}^\perp = \left\{w \in \mathbb{F}^n_2 \mid \forall v \in \mathcal{C}, \langle w, v \rangle = 0\right\}
.\end{equation}
The orthogonal complement has dimension $\dim(\mathcal{C}^\perp) = n - k$. Thus, $\mathcal{C}^\perp$ has a basis formed by $n - k$ linearly independent vectors: $\{h_1, h_2, \ldots, h_{n - k}\}$. We construct the parity-check matrix $H$ by setting these basis vectors as its rows:
\begin{equation}
H = \begin{pmatrix} - & h_1 & - \\ - & h_2 & - \\ & \vdots & \\ - & h_{n - k} & - \end{pmatrix}
.\end{equation}

\begin{itemize}
\item ($\implies$) Let $v \in \mathcal{C}$. By the definition of $\mathcal{C}^\perp$, $v$ is orthogonal to every basis vector $h_i$. Therefore, the dot product $h_i \cdot v = 0$ for all $i$, which yields $H v = 0$.
\item ($\impliedby$) Let $v \in \mathbb{F}^n_2$ such that $H v = 0$. This means that $v$ is orthogonal to every row $h_i$ of $H$. Since the rows of $H$ span $\mathcal{C}^\perp$, $v$ is orthogonal to all of $\mathcal{C}^\perp$, implying that $v \in (\mathcal{C}^\perp)^\perp$. By Lemma~\ref{lemma:complement_of_orthogonal_complement}, we conclude that $v \in \mathcal{C}$.
\end{itemize}
\end{proof}
\end{theorem}

The existence of a parity-check matrix provides a systematic framework for characterizing linear codes. During the detection phase, checking whether a received vector contains an error reduces to a simple matrix multiplication: applying the parity-check matrix to the received vector to compute the error syndrome. This process is best illustrated with a concrete example:

\begin{example}\label{ex:repetition_code}
The 3-bit repetition code is a simple $[3, 1]$ linear code with generator matrix:
\begin{equation}
G = \begin{bmatrix}1 \\ 1 \\ 1\end{bmatrix}
.\end{equation}
It encodes a single bit of information $u \in \mathbb{F}_2$ into a 3-bit codeword $v = Gu$:
\begin{equation}
\begin{bmatrix}1 \\ 1 \\ 1\end{bmatrix} \begin{bmatrix} 0 \end{bmatrix} = \begin{bmatrix} 0 \\ 0 \\ 0 \end{bmatrix};\quad \begin{bmatrix}1 \\ 1 \\ 1\end{bmatrix} \begin{bmatrix} 1 \end{bmatrix} = \begin{bmatrix} 1 \\ 1 \\ 1 \end{bmatrix}
.\end{equation}

Its corresponding parity-check matrix $H$ must satisfy $H G = 0$. A valid choice is:
\begin{equation}
H = \begin{bmatrix}1 & 1 & 0 \\ 1 & 0 & 1\end{bmatrix}
.\end{equation}

When a vector $r \in \mathbb{F}_2^3$ is received, we can check for errors by calculating the \textit{syndrome} vector $s = H r$. If $s = 0$, the received vector is a valid codeword. If a single-bit error occurred, the syndrome $s$ matches one of the columns of $H$:

\begin{itemize}
\item If $s = \begin{bmatrix} 0 \\ 0 \end{bmatrix}$, there is no error (or an undetectable error occurred).
\item If $s = \begin{bmatrix} 1 \\ 1 \end{bmatrix}$, error occurred in the first bit (matches the 1st column of $H$).
\item If $s = \begin{bmatrix} 1 \\ 0 \end{bmatrix}$, error occurred in the second bit (matches the 2nd column of $H$).
\item If $s = \begin{bmatrix} 0 \\ 1 \end{bmatrix}$, error occurred in the third bit (matches the 3rd column of $H$).
\end{itemize}
\end{example}
